\documentclass[11pt]{article}

\usepackage[T1]{fontenc}
\usepackage[utf8]{inputenc}
\usepackage{amsmath,amssymb,amsfonts,amsthm}
\usepackage{bm}
\usepackage{booktabs}
\usepackage[margin=1in]{geometry}
\usepackage[hidelinks]{hyperref}
\usepackage{framed}

% --- notation ---------------------------------------------------------------
\newcommand{\R}{\mathbb{R}}
\newcommand{\Ncal}{\mathcal{N}}
\newcommand{\x}{\bm{x}}
\newcommand{\z}{\bm{z}}
\newcommand{\uu}{\bm{u}}
\newcommand{\yy}{\bm{y}}
\newcommand{\eps}{\bm{\varepsilon}}
\newcommand{\Fhat}{\widehat{F}}
\newcommand{\fhat}{\widehat{f}}
\newcommand{\Qhat}{\widehat{Q}}
\DeclareMathOperator{\sd}{sd}
\DeclareMathOperator{\Var}{Var}

\newtheorem*{deriv}{Derivation}
\theoremstyle{remark}
\newtheorem*{plain}{In words}

\title{\textbf{The coppuccino copula flow, derived from scratch}\\[2pt]
\large A plain--language note with all the mathematics}
\author{}
\date{}

\begin{document}
\maketitle

\begin{framed}
\noindent
\textbf{Status: tutorial derivation of the 1.0.0 marginal, not the shipped one.}
These notes derive the copula--flow construction using the
\emph{monotone piecewise--cubic (PCHIP) double--spline} marginal, which is the
model coppuccino shipped in 1.0.0. coppuccino~2.0.0 removed it, together with
the Gaussian tail extension of Sec.~\ref{sec:tail} (\texttt{tail\_extension=True}),
so the marginal sections (Secs.~\ref{sec:pchip}--\ref{sec:tail}) describe code
that no longer exists. What remains accurate:
Sklar's theorem and the copula factorization (Sec.~\ref{sec:copula}), the
change--of--variables identity and the marginal Jacobian
(Secs.~\ref{sec:problem}, \ref{sec:gaussianization}), monotone Hermite
interpolation as such (Sec.~\ref{sec:pchip}; the shipped marginal still uses
PCHIP slopes at its knots), the triangular spline copula flow
(Sec.~\ref{sec:copula-flow}), and the assembled model density and sampling
procedure (Secs.~\ref{sec:master}, \ref{sec:algorithms}).

\medskip\noindent
What they do \emph{not} describe: the only marginal since 2.0.0, the
smooth--tail map --- a single clip--free monotone map of all of
$\R$ onto the support (a rational--quadratic--spline bulk at the data quantiles
joined to analytic tails), whose forward map, inverse and derivative share one
parameterization rather than being two independently fitted splines. That
family is what the boundary discussion of Sec.~\ref{sec:why} should be read
against; the caveat stated there is the PCHIP--era one. The notes were written
against coppuccino~1.x; for 2.0.0 only this box was revised.
\texttt{coppuccino\_notes.pdf} in this directory is a build
artifact of an \emph{earlier} revision of this source: it predates this box and
is not rebuilt automatically, so where the two disagree the \texttt{.tex} is
authoritative. Rebuild it with
\texttt{pdflatex coppuccino\_notes.tex} run three times (once for the aux file,
once for the table of contents, once for the cross--references).
\end{framed}

\begin{abstract}
\noindent
This is a companion to the terse methods description: same architecture, but
nothing assumed. We build the model one piece at a time and derive every formula,
and we stop to explain the machine--learning vocabulary (a ``normalizing flow'',
``training'', a ``base distribution'', why the Jacobian is cheap) in ordinary
terms. If you are comfortable with calculus and probability but have never taken
a machine--learning course, this is written for you. The target is the
Gaussianized--marginal pipeline with monotone piecewise--cubic (PCHIP)
marginals --- see the status box above for how that relates to what the library
fits today.
\end{abstract}

\tableofcontents

\section{The problem, and the one trick we use everywhere}
\label{sec:problem}

We have $N$ samples $\{\x^{(n)}\}_{n=1}^{N}$ drawn from some unknown continuous
density $p(\x)$ on $\R^{d}$ --- think of the posterior draws that come out of an
MCMC sampler. We want two things: to \emph{evaluate} $p(\x)$ at new points, and
to \emph{draw} fresh samples from it. A model that does both lets us resample a
posterior cheaply, compute importance weights, and run calibration checks.

Everything below rests on a single identity from calculus: the
\emph{change--of--variables} formula for densities. It is worth deriving once,
carefully, because we will use it three times.

\subsection{Change of variables for densities}

Let $g:\R^{d}\to\R^{d}$ be a smooth, invertible map (a \emph{diffeomorphism}),
and let $\yy=g(\x)$. If $\x$ has density $p_X$, what is the density $p_Y$ of
$\yy$?

\begin{deriv}
Probability mass is conserved: for any region $A\subseteq\R^{d}$, the chance that
$\x\in A$ must equal the chance that $\yy\in g(A)$,
\begin{equation}
\int_{A} p_X(\x)\,\mathrm{d}\x \;=\; \int_{g(A)} p_Y(\yy)\,\mathrm{d}\yy .
\end{equation}
Change the integration variable on the right from $\yy$ back to $\x$ via
$\yy=g(\x)$. The volume element transforms by the absolute Jacobian determinant,
$\mathrm{d}\yy=\bigl|\det J_g(\x)\bigr|\,\mathrm{d}\x$, where
$J_g(\x)=\partial g/\partial\x$ is the $d\times d$ matrix of partial derivatives.
Substituting,
\begin{equation}
\int_{A} p_X(\x)\,\mathrm{d}\x
= \int_{A} p_Y\bigl(g(\x)\bigr)\,\bigl|\det J_g(\x)\bigr|\,\mathrm{d}\x .
\end{equation}
This holds for \emph{every} region $A$, so the integrands must agree pointwise,
\begin{equation}
\label{eq:cov-forward}
p_X(\x) = p_Y\bigl(g(\x)\bigr)\,\bigl|\det J_g(\x)\bigr| .
\end{equation}
Solving for $p_Y$ (write $\x=g^{-1}(\yy)$ and use that the Jacobian of the
inverse is the inverse of the Jacobian, so its determinant is the reciprocal),
\begin{equation}
\label{eq:cov}
\boxed{\;
p_Y(\yy) = p_X\bigl(g^{-1}(\yy)\bigr)\,
\Bigl|\det J_{g^{-1}}(\yy)\Bigr|
= \frac{p_X\bigl(g^{-1}(\yy)\bigr)}{\bigl|\det J_g(g^{-1}(\yy))\bigr|}. \;}
\end{equation}
\end{deriv}

\begin{plain}
A density is ``probability per unit volume''. When a map stretches a region, the
same probability is spread over a larger volume, so the density drops --- by
exactly the local volume--stretch factor, which is $|\det J|$. That factor is the
only thing the change of variables adds. In one dimension it is just
$|\mathrm{d}y/\mathrm{d}x|$; the determinant is its multivariate generalization.
\end{plain}

Taking logarithms of \eqref{eq:cov} turns the product into a sum, which is how it
always appears below:
\begin{equation}
\label{eq:logcov}
\log p_Y(\yy) = \log p_X\bigl(g^{-1}(\yy)\bigr) + \log\Bigl|\det J_{g^{-1}}(\yy)\Bigr| .
\end{equation}
Two facts we will lean on. (i)~If $g$ acts on each coordinate separately,
$y_i=g_i(x_i)$, then $J_g$ is \emph{diagonal} and
$\log|\det J_g|=\sum_i\log|g_i'(x_i)|$ --- a sum of one--dimensional terms.
(ii)~If $g$ is \emph{triangular} ($y_i$ depends only on $x_1,\dots,x_i$), then
$J_g$ is triangular and $\log|\det J_g|=\sum_i\log|\partial y_i/\partial x_i|$ ---
still just the diagonal. Both cases cost $O(d)$ instead of the $O(d^3)$ of a
general determinant. This is the entire reason the architecture is built from
coordinatewise and triangular maps.

\section{What a normalizing flow is}
\label{sec:flow-intro}

A \emph{normalizing flow} is a way to build a flexible density that you can both
evaluate and sample from. The recipe has two ingredients:
\begin{enumerate}
\item a \emph{base distribution} --- a simple density you can already sample and
      evaluate, almost always a standard Gaussian $\eps\sim\Ncal(\bm0,I_d)$;
\item a \emph{learnable invertible map} $T_\theta:\R^d\to\R^d$, where $\theta$ is
      a vector of adjustable numbers (the ``parameters'').
\end{enumerate}
The model \emph{samples} by drawing $\eps$ from the base and pushing it through
the map, $\x=T_\theta(\eps)$. Its density follows immediately from the change of
variables \eqref{eq:cov}, with the base playing the role of $p_X$ and $T_\theta$
the role of $g^{-1}$'s inverse:
\begin{equation}
\label{eq:flow-density}
\log p_\theta(\x)
= \underbrace{\log\Ncal\!\bigl(T_\theta^{-1}(\x);\,\bm0,I_d\bigr)}_{\text{base density at the pre-image}}
\;+\;
\underbrace{\log\Bigl|\det \tfrac{\partial T_\theta^{-1}}{\partial\x}(\x)\Bigr|}_{\text{volume correction}} .
\end{equation}

\begin{plain}
The name says what it does: run \emph{backwards} (data $\to$ base), a good flow
turns the complicated data distribution into the simple, ``normalized'' Gaussian.
Run forwards, it turns Gaussian noise into data. Because the map is invertible
and we can compute its Jacobian, \eqref{eq:flow-density} gives an exact density
--- not an approximation or a bound.
\end{plain}

\paragraph{What ``learnable'' and ``training'' mean.}
The parameters $\theta$ start at some arbitrary value and are adjusted until the
model fits the data. ``Fits'' is made precise by \emph{maximum likelihood}: we
want the density that makes the observed samples as probable as possible, i.e.\
we maximize the average log--density of the data,
\begin{equation}
\label{eq:mle}
\theta^\star = \arg\max_\theta\ \frac1N\sum_{n=1}^N \log p_\theta\bigl(\x^{(n)}\bigr).
\end{equation}
Maximizing \eqref{eq:mle} is the same as minimizing the Kullback--Leibler
divergence from the data distribution to the model, so the fitted model is the
closest reachable density to the truth in that sense. The maximization is done by
\emph{gradient ascent}: compute the gradient of \eqref{eq:mle} with respect to
$\theta$ (automatic differentiation does this for us), take a small step uphill,
repeat. ``Adam'' is just a popular gradient--ascent rule with per--parameter
adaptive step sizes; ``epoch'' means one pass through the data; ``early
stopping'' means we watch the objective on a held--out slice of the samples and
stop when it stops improving, to avoid overfitting. None of this changes the
model --- it only chooses $\theta$.

\paragraph{How $T_\theta$ is built, and why it stays invertible.}
A single simple invertible map is not expressive enough, so $T_\theta$ is a
\emph{composition} of many simple invertible layers,
$T_\theta=T_L\circ\cdots\circ T_1$. A composition of invertible maps is
invertible, and by the chain rule its log--Jacobian is the sum of the layers',
so we only ever need each layer to be invertible with a cheap Jacobian. Two layer
types, alternated, make up the flow used here (Section~\ref{sec:copula-flow});
both were chosen precisely so their Jacobians are diagonal or triangular, per the
remark at the end of Section~\ref{sec:problem}. (Many flows in the literature use
small neural networks inside the layers; the one here does not --- it is built
from monotone splines and triangular linear maps, which keeps it simple to
reason about.)

\section{The copula idea: separate the marginals from the dependence}
\label{sec:copula}

Fitting a flow directly to raw data is hard when the individual coordinates have
awkward shapes --- sharp peaks, hard boundaries, heavy tails, multiple modes.
The flow burns its capacity reshaping each axis instead of learning how the axes
relate. The copula construction removes that burden.

\subsection{Sklar's theorem, and why it is almost obvious}

Write the marginal CDF of coordinate $i$ as $F_i$ and its density as
$p_i=F_i'$. Define the \emph{probability integral transform} (PIT) of each
coordinate, $u_i=F_i(x_i)$.

\begin{deriv}[the PIT makes a coordinate uniform]
If $X$ has continuous, strictly increasing CDF $F$, then $U=F(X)$ is uniform on
$(0,1)$. Compute its CDF directly, for $u\in(0,1)$:
\begin{equation}
\Pr(U\le u)=\Pr\bigl(F(X)\le u\bigr)=\Pr\bigl(X\le F^{-1}(u)\bigr)
= F\bigl(F^{-1}(u)\bigr)=u,
\end{equation}
which is the CDF of a $\mathrm{Uniform}(0,1)$ variable. (We used that $F$ is
increasing to invert the inequality.)
\end{deriv}

So after the PIT, every coordinate of $\uu=(u_1,\dots,u_d)$ is marginally
uniform, while the \emph{joint} arrangement of $\uu$ still records how the
original coordinates move together. That joint law is the \emph{copula}: by
definition, a copula $C$ is the joint CDF of $\uu$, and its density $c$ carries
all of the dependence and none of the marginal shape (each marginal of $c$ is
flat). Sklar's theorem is the statement that this always factors,
\begin{equation}
\label{eq:sklar}
p(\x)=c\bigl(F_1(x_1),\dots,F_d(x_d)\bigr)\ \prod_{i=1}^{d} p_i(x_i).
\end{equation}

\begin{deriv}[Sklar's density, from the change of variables]
The map $\x\mapsto\uu$ with $u_i=F_i(x_i)$ is coordinatewise, so by
\eqref{eq:cov-forward} and fact~(i),
$p(\x)=c(\uu)\prod_i|\mathrm{d}u_i/\mathrm{d}x_i|
      =c(\uu)\prod_i F_i'(x_i)=c(\uu)\prod_i p_i(x_i)$,
which is exactly \eqref{eq:sklar}. So Sklar's theorem is just the change of
variables applied to the PIT; the density of $\uu$ \emph{is} the copula density
$c$.
\end{deriv}

The plan writes itself: estimate the $d$ marginals $F_i$ nonparametrically (they
absorb every awkward one--dimensional shape), and model the copula $c$ with a
flow. The flow now only has to represent dependence, which is a much easier job.

\subsection{Why go through a Gaussian instead of staying on the cube}

One could model $c$ directly on the unit cube $[0,1]^d$. The pipeline described
here instead sends each uniform coordinate one step further, to a standard normal,
\begin{equation}
\label{eq:gaussianize}
z_i=\Phi^{-1}(u_i)=\Phi^{-1}\bigl(F_i(x_i)\bigr),
\end{equation}
where $\Phi$ is the standard--normal CDF (so $\Phi^{-1}$ is its quantile
function). By the same PIT argument in reverse, $z_i$ is then marginally standard
normal. The advantage is that the flow now lives on unbounded space with a
Gaussian base that already matches the marginals, so it need not fight the hard
walls of the cube; the price, discussed at the end, is boundary behavior. We call
$M:\x\mapsto\z$ the \emph{marginal Gaussianization map}. It acts on each
coordinate independently.

In the library both choices are still built. The single--flow entry point
\texttt{normalizing\_flows\_fit} is the Gaussianized pipeline derived here; the
many--source entry point \texttt{normalizing\_flows\_fit\_fleet} models the
copula on the cube directly, over a uniform base, and is the one to reach for
when the marginals run flat into a hard wall.

\section{Architecture overview}
\label{sec:arch}

Two stages, one frozen and one trained. Reading the pipeline forwards
(generative) and backwards (evaluation):
\begin{equation}
\label{eq:pipeline}
\eps\;\sim\;\Ncal(\bm0,I_d)
\ \xrightarrow[\text{copula flow (trained)}]{\ T_\theta\ }\ \z
\ \xrightarrow[\text{Gaussianization (frozen)}]{\ M^{-1}\ }\ \x .
\end{equation}
Stage~1 is the coordinatewise map $M$ of \eqref{eq:gaussianize}, fixed once the
marginals $\Fhat_i$ are estimated. Stage~2 is the flow $T_\theta$ of
Section~\ref{sec:flow-intro}, which models the joint law $p_Z$ of the Gaussianized
data. The next three sections derive each stage and then compose them.

\section{Stage 1: the marginal Gaussianization, differentiated}
\label{sec:gaussianization}

We need the Jacobian of $M$ to use the change of variables. Because $M$ is
coordinatewise, we only need one coordinate's derivative; drop the index $i$.
The map is $z=\Phi^{-1}(F(x))$, a composition, so use the chain rule with
$u=F(x)$:
\begin{equation}
\frac{\mathrm{d}z}{\mathrm{d}x}
=\frac{\mathrm{d}z}{\mathrm{d}u}\,\frac{\mathrm{d}u}{\mathrm{d}x}.
\end{equation}
The second factor is $\mathrm{d}u/\mathrm{d}x=F'(x)=f(x)$, the marginal density.
For the first factor we differentiate an inverse function.

\begin{deriv}[derivative of $\Phi^{-1}$]
Let $\phi=\Phi'$ be the standard--normal density. From $\Phi(\Phi^{-1}(u))=u$,
differentiate both sides in $u$: $\phi(\Phi^{-1}(u))\cdot(\Phi^{-1})'(u)=1$, so
\begin{equation}
(\Phi^{-1})'(u)=\frac{1}{\phi(\Phi^{-1}(u))}=\frac{1}{\phi(z)} .
\end{equation}
\end{deriv}

\noindent Combining,
\begin{equation}
\label{eq:marginal-jac}
\boxed{\;\frac{\mathrm{d}z}{\mathrm{d}x}=\frac{f(x)}{\phi(z)},
\qquad z=\Phi^{-1}\bigl(F(x)\bigr).\;}
\end{equation}
Restoring the coordinate index and using fact~(i), the log--Jacobian of the full
map $M$ is the sum
\begin{equation}
\label{eq:M-logdet}
\log\Bigl|\det \tfrac{\partial M}{\partial\x}\Bigr|
=\sum_{i=1}^{d}\Bigl[\log f_i(x_i)-\log\phi(z_i)\Bigr].
\end{equation}
The inverse map, used when we sample, undoes each step: apply the normal CDF, then
the marginal quantile function,
\begin{equation}
\label{eq:M-inverse}
x_i=F_i^{-1}\bigl(\Phi(z_i)\bigr).
\end{equation}
Everything in \eqref{eq:marginal-jac}--\eqref{eq:M-inverse} needs three
one--dimensional objects per coordinate: the CDF $F$, its density $f=F'$, and its
inverse $F^{-1}$. Building those from the samples is the next section.

\section{Stage 1, continued: estimating the marginal with a monotone spline}
\label{sec:pchip}

We estimate each marginal CDF as a smooth, strictly increasing curve through the
empirical CDF, using a monotone piecewise--cubic Hermite interpolant (PCHIP).
Drop the coordinate index.

\subsection{Knots and the empirical CDF}

Choose $K$ knots at empirical quantiles of the sample,
\begin{equation}
\label{eq:knots}
x_g^{(m)}=\Qhat\!\Bigl(\tfrac{m}{K-1}\Bigr),\qquad m=0,\dots,K-1,
\end{equation}
where $\Qhat$ is the sample quantile function. Spacing the knots evenly in
\emph{probability} (not in $x$) places them densely where the data are dense.
The empirical CDF value at a knot is the fraction of samples at or below it,
\begin{equation}
\label{eq:ecdf}
\Fhat^{(m)}=\frac1N\,\#\bigl\{n:x^{(n)}\le x_g^{(m)}\bigr\}.
\end{equation}
The knot count is set adaptively,
$K=\max\!\bigl(20,\min(K_{\max},\lfloor N/3\rfloor)\bigr)$ with default
$K_{\max}=200$, so small samples do not get an over--fine grid.

\subsection{Cubic Hermite interpolation, and the basis}

Between two adjacent knots we want a cubic that matches the CDF value \emph{and}
a chosen slope at each end (matching slopes is what makes the joined curve
smooth, i.e.\ $C^1$). On $[x_g^{(m)},x_g^{(m+1)}]$ set
$h_m=x_g^{(m+1)}-x_g^{(m)}$ and the local coordinate $t=(x-x_g^{(m)})/h_m\in[0,1]$.
Writing the endpoint values $y_0=\Fhat^{(m)},y_1=\Fhat^{(m+1)}$ and the (rescaled)
endpoint slopes, the interpolant is
\begin{equation}
\label{eq:hermite}
\Fhat(x)=\Fhat^{(m)}H_{00}(t)+h_m d_m H_{10}(t)
        +\Fhat^{(m+1)}H_{01}(t)+h_m d_{m+1}H_{11}(t),
\end{equation}
with the cubic Hermite basis
\begin{equation}
\label{eq:hbasis}
\begin{aligned}
H_{00}(t)&=2t^3-3t^2+1, & H_{10}(t)&=t^3-2t^2+t,\\
H_{01}(t)&=-2t^3+3t^2,  & H_{11}(t)&=t^3-t^2.
\end{aligned}
\end{equation}

\begin{deriv}[where the basis comes from]
Each basis cubic is the unique degree--3 polynomial isolating one endpoint
condition. Take $H_{00}$: we require value $1$ at $t=0$, value $0$ at $t=1$, and
slope $0$ at both ends. Writing $H_{00}=a+bt+ct^2+et^3$, the four conditions
$H_{00}(0)=1,\ H_{00}'(0)=0,\ H_{00}(1)=0,\ H_{00}'(1)=0$ give $a=1,\ b=0$, and
$c+e=-1,\ 2c+3e=0$, hence $e=2,c=-3$, i.e.\ $H_{00}=1-3t^2+2t^3$. The other three
follow the same way. One checks that
$\Fhat(x_g^{(m)})=\Fhat^{(m)}$, $\Fhat(x_g^{(m+1)})=\Fhat^{(m+1)}$,
$\Fhat'(x_g^{(m)})=d_m$, and $\Fhat'(x_g^{(m+1)})=d_{m+1}$ --- the factor $h_m$ in
\eqref{eq:hermite} converts the slope in $x$ to a slope in $t$ (chain rule).
\end{deriv}

\subsection{Choosing the slopes so the CDF cannot wiggle}

Nothing so far forces $\Fhat$ to be increasing --- and a CDF that dips would give
a negative density, which is nonsense. The slopes $d_m$ must be chosen to keep the
interpolant monotone. Let $\delta_m=(\Fhat^{(m+1)}-\Fhat^{(m)})/h_m\ge0$ be the
secant slope across interval $m$.

\begin{plain}
Intuition: if the data say the CDF rises gently across one bin and steeply across
the next, the knot between them should have a slope somewhere in between --- and
if the data flatten out (a local plateau), the slope there should be zero so the
curve does not overshoot and come back down.
\end{plain}

Fritsch and Carlson \cite{FritschCarlson1980} made this precise: the Hermite
piece is monotone on a bin whenever the endpoint slopes lie in a bounded region
around the secant (a sufficient version is $0\le d_m\le 3\min(\delta_{m-1},
\delta_m)$). PCHIP picks the slopes to sit safely inside that region: a knot at a
local extremum ($\delta_{m-1}\delta_m\le0$) gets $d_m=0$, and otherwise $d_m$ is
the \emph{weighted harmonic mean} of the two neighbouring secants,
\begin{equation}
\label{eq:fc}
d_m=\frac{w_1+w_2}{\,w_1/\delta_{m-1}+w_2/\delta_m\,},
\qquad w_1=2h_m+h_{m-1},\quad w_2=h_m+2h_{m-1},
\end{equation}
with a one--sided formula at the two ends of the grid.

\begin{plain}
Why a harmonic mean rather than an ordinary average? The harmonic mean of two
positive numbers is dominated by the \emph{smaller} one and collapses to zero if
either does. That is exactly the safe behaviour: near a flat bin the slope is
pulled toward zero, which is what prevents overshoot. The result is guaranteed
nondecreasing, so $\Fhat$ is a genuine CDF.
\end{plain}

\subsection{The density and the quantile function for free}

The marginal density is just the analytic derivative of \eqref{eq:hermite} --- a
continuous, nonnegative, piecewise--quadratic function,
\begin{equation}
\label{eq:pdf}
\fhat(x)=\Fhat'(x)=\frac{1}{h_m}\Bigl[\Fhat^{(m)}H_{00}'(t)+h_m d_m H_{10}'(t)
+\Fhat^{(m+1)}H_{01}'(t)+h_m d_{m+1}H_{11}'(t)\Bigr],
\end{equation}
(the $1/h_m$ is again the chain rule, $\mathrm{d}t/\mathrm{d}x$). Having $\fhat$
in closed form is what makes the marginal Jacobian \eqref{eq:marginal-jac} cheap
--- no numerical differentiation. The quantile function $F^{-1}$ needed for
sampling is a \emph{second} PCHIP interpolant, this one through the swapped knots
$\{(\Fhat^{(m)},x_g^{(m)})\}$ (CDF value on the horizontal axis, $x$ on the
vertical). Because the two splines are fit independently, they invert each other
only up to the interpolation error, which shrinks as the knot count $K$ grows
(and saturates at $K_{\max}$).

\subsection{Beyond the data range: the Gaussian tail}
\label{sec:tail}

The knots \eqref{eq:knots} only cover $[x_{\min},x_{\max}]$. By default the CDF
is held flat at its boundary value outside that range (appropriate when a prior
bounds the parameter). To extrapolate more gracefully on unbounded parameters,
continue the CDF with a Gaussian tail matched at the boundary. Fit a Gaussian to
the coordinate, $\mu=\bar x$, $\sigma=\sd(x)$, with CDF
$\Phi_{\mu,\sigma}(x)=\Phi((x-\mu)/\sigma)$, and require the tail to meet the
interpolant continuously.

\begin{deriv}[the matching constants]
On the left we want a scaled Gaussian CDF, $\Fhat(x)=c_{\mathrm{lo}}\,
\Phi_{\mu,\sigma}(x)$ for $x<x_{\min}$, that agrees with the interpolant's
boundary value $F_{\mathrm{lo}}=\Fhat(x_{\min})$. Setting $x=x_{\min}$ and solving,
\begin{equation}
c_{\mathrm{lo}}\,\Phi_{\mu,\sigma}(x_{\min})=F_{\mathrm{lo}}
\ \Longrightarrow\
c_{\mathrm{lo}}=\frac{F_{\mathrm{lo}}}{\Phi_{\mu,\sigma}(x_{\min})}.
\end{equation}
On the right we scale the \emph{survival} function so that
$\Fhat(x)=1-c_{\mathrm{hi}}\bigl[1-\Phi_{\mu,\sigma}(x)\bigr]$ matches
$F_{\mathrm{hi}}=\Fhat(x_{\max})$ at $x=x_{\max}$:
\begin{equation}
1-c_{\mathrm{hi}}\bigl[1-\Phi_{\mu,\sigma}(x_{\max})\bigr]=F_{\mathrm{hi}}
\ \Longrightarrow\
c_{\mathrm{hi}}=\frac{1-F_{\mathrm{hi}}}{1-\Phi_{\mu,\sigma}(x_{\max})}.
\end{equation}
\end{deriv}

\noindent Putting the three pieces together,
\begin{equation}
\label{eq:tail-cdf}
\Fhat(x)=
\begin{cases}
c_{\mathrm{lo}}\,\Phi_{\mu,\sigma}(x), & x<x_{\min},\\[3pt]
\text{PCHIP interpolant \eqref{eq:hermite}}, & x_{\min}\le x\le x_{\max},\\[3pt]
1-c_{\mathrm{hi}}\bigl[1-\Phi_{\mu,\sigma}(x)\bigr], & x>x_{\max}.
\end{cases}
\end{equation}
The tail density is the derivative, $c_{\mathrm{lo}}\phi_{\mu,\sigma}(x)$ on the
left and $c_{\mathrm{hi}}\phi_{\mu,\sigma}(x)$ on the right, and the tail quantile
inverts in closed form, e.g.\ on the left
$x=\mu+\sigma\,\Phi^{-1}(u/c_{\mathrm{lo}})$. The match is continuous in the CDF;
the density can have a kink at the join, which is harmless for the transform.

\section{Stage 2: the copula flow}
\label{sec:copula-flow}

The flow $T_\theta$ models the joint density $p_Z$ of the Gaussianized vector
$\z$. Concretely it is a stack of $L$ layers, each combining two invertible
pieces whose Jacobians we already know how to handle cheaply.

\paragraph{Piece one: an elementwise monotone spline.}
Each coordinate is passed through its own increasing rational--quadratic spline
--- a curve stitched from pieces that are each a ratio of two quadratics,
constrained to be monotone with adjustable knot heights, widths, and slopes.
These knot parameters live in $\theta$ and are what training adjusts. Because the
spline acts coordinate by coordinate, its Jacobian is diagonal and its
log--determinant is $\sum_i\log(\text{spline slope})_i$. (A $\tanh$ squashes the
input to a finite window first, so the spline only has to be defined on a bounded
interval; the constant \texttt{tanh\_max\_val} sets that window.) This piece adds
nonlinearity but, on its own, cannot make coordinates interact.

\paragraph{Piece two: a triangular linear map.}
The coordinates are then mixed by an affine map $\z\mapsto L\z+\bm b$ with $L$ a
triangular matrix. Triangularity is the whole point: by fact~(ii) of
Section~\ref{sec:problem} the Jacobian is $L$, its determinant is the product of
the diagonal, and $\log|\det|=\sum_i\log|L_{ii}|$. This is what lets one
coordinate's transformation depend on the others --- i.e.\ what actually models
dependence --- without an expensive determinant. Alternating the orientation of
$L$ (lower, then upper, \dots) across layers lets every coordinate eventually
influence every other.

\paragraph{The flow's density.}
Stacking these layers gives an invertible $T_\theta$ whose log--Jacobian is the
sum of the per--layer diagonal terms, so \eqref{eq:flow-density} specialises to
\begin{equation}
\label{eq:pz}
\log p_Z(\z)=\log\Ncal\!\bigl(T_\theta^{-1}(\z);\bm0,I_d\bigr)
+\sum_{\ell=1}^{L}\log\Bigl|\det \tfrac{\partial T_\ell^{-1}}{\partial(\cdot)}\Bigr|,
\end{equation}
every term of which is an $O(d)$ sum over a diagonal. Because the marginals were
already standardized in Stage~1, the base $\Ncal(\bm0,I_d)$ is a good starting
guess for $\z$, and the flow only has to bend it into the right dependence.

\paragraph{Training.}
The parameters $\theta$ are fit by the maximum likelihood of
Section~\ref{sec:flow-intro}. Since the marginal maps are frozen, the marginal
term in the full density (next section) does not depend on $\theta$, so
maximizing the data likelihood is equivalent to maximizing
$\sum_n\log p_Z\bigl(M(\x^{(n)})\bigr)$ on the fixed Gaussianized data --- we
literally transform the data once, then train an ordinary flow on it (Adam, early
stopping on a held--out split).

\section{Putting the stages together: the model density}
\label{sec:master}

Now compose. The model is $\x=M^{-1}(T_\theta(\eps))$, so its density is the flow
density for $\z$ pushed through the marginal map $M$. Apply the change of
variables \eqref{eq:logcov} to $M$, using its log--Jacobian \eqref{eq:M-logdet}:
\begin{equation}
\label{eq:master}
\boxed{\;
\log p(\x)=\log p_Z(\z)+\sum_{i=1}^{d}\Bigl[\log \fhat_i(x_i)-\log\phi(z_i)\Bigr],
\qquad z_i=\Phi^{-1}\bigl(\Fhat_i(x_i)\bigr).\;}
\end{equation}
The first term is the flow \eqref{eq:pz}; the sum is the marginal correction. That
is the whole model density.

\begin{deriv}[the same formula, read as a copula --- a good consistency check]
Sklar \eqref{eq:sklar} says $\log p(\x)=\log c(\uu)+\sum_i\log\fhat_i(x_i)$ with
$u_i=\Fhat_i(x_i)$. We only need to see that the flow supplies $\log c$. The
copula density is the density of $\uu$; since $\z=\Phi^{-1}(\uu)$ coordinatewise,
the change of variables \eqref{eq:cov-forward} with fact~(i) gives
\begin{equation}
c(\uu)=p_Z(\z)\prod_{i}\Bigl|\frac{\mathrm{d}z_i}{\mathrm{d}u_i}\Bigr|
      =p_Z(\z)\prod_i\frac{1}{\phi(z_i)},
\quad\text{i.e.}\quad
\log c(\uu)=\log p_Z(\z)-\sum_i\log\phi(z_i),
\end{equation}
using $\mathrm{d}z_i/\mathrm{d}u_i=(\Phi^{-1})'(u_i)=1/\phi(z_i)$. Substituting
into Sklar reproduces \eqref{eq:master} exactly. So ``the flow's density minus the
Gaussian marginal terms'' \emph{is} the log--copula --- the two viewpoints agree.
\end{deriv}

\section{Sampling and evaluation, as procedures}
\label{sec:algorithms}

Both directions of the pipeline \eqref{eq:pipeline} are now explicit.

\medskip
\noindent\textbf{To evaluate} $\log p(\x)$ at a point:
\begin{enumerate}\itemsep2pt
\item Gaussianize: $u_i\leftarrow\Fhat_i(x_i)$, then $z_i\leftarrow\Phi^{-1}(u_i)$, for each $i$.
\item Return $\log p_Z(\z)+\sum_i\bigl[\log\fhat_i(x_i)-\log\phi(z_i)\bigr]$.
\end{enumerate}

\noindent\textbf{To sample} $\x\sim p$:
\begin{enumerate}\itemsep2pt
\item Draw $\eps\sim\Ncal(\bm0,I_d)$ and push it through the flow, $\z\leftarrow T_\theta(\eps)$.
\item Invert the marginals, $x_i\leftarrow\Fhat_i^{-1}\bigl(\Phi(z_i)\bigr)$, for each $i$.
\end{enumerate}
Doing both in one forward pass returns a sample together with its exact
log--density, which is what importance weighting needs.

\section{Why the pieces were chosen, and one caveat}
\label{sec:why}

\begin{itemize}
\item \textbf{Separation of concerns.} The nonparametric marginals \eqref{eq:hermite}
absorb every hard one--dimensional shape, so the flow \eqref{eq:pz} spends all its
parameters on dependence. Each half does the job it is good at.
\item \textbf{Everything is closed--form.} The marginal CDF, its density, and (up
to interpolation error) its inverse all come from cheap spline evaluations; the
flow's density and inverse are exact. So \eqref{eq:master} is an exact density,
not a bound, and samples come with honest weights.
\item \textbf{Cheap Jacobians throughout.} Every map in the pipeline is
coordinatewise or triangular, so every determinant is an $O(d)$ sum over a
diagonal --- the reason the whole thing scales.
\item \textbf{Caveat at hard boundaries.} Gaussianization sends the ends of the
range to $\pm\infty$ through $\Phi^{-1}$, so a query \emph{at} a hard bound
arrives at the copula flow at a very large $|z|$, where the flow is
extrapolating past its outermost knot. The density it reports there is
self--consistent with the sampler --- draws and \texttt{log\_prob} come from one
parameterization --- but it is not a trustworthy pointwise value, and on a chain
that really has mass sitting on a wall it is much \emph{larger} than the
interior density rather than smaller. The Gaussian pipeline is the right choice
on unbounded parameters and whenever the tail models of
Section~\ref{sec:tail} are needed; when the density at the wall itself is the
quantity of interest, the cube--native fleet pipeline never leaves $[0,1]^d$ and
does not have this failure mode. Note also that the double--spline construction
of Section~\ref{sec:pchip} inverts only up to interpolation error; the
smooth--tail marginal named in the status box replaces it with a
single parameterization whose inverse is exact.
\end{itemize}

\begin{thebibliography}{9}

\bibitem{Sklar1959}
A.~Sklar,
\emph{Fonctions de r\'epartition \`a $n$ dimensions et leurs marges},
Publ.\ Inst.\ Statist.\ Univ.\ Paris \textbf{8}, 229--231 (1959).

\bibitem{FritschCarlson1980}
F.~N.~Fritsch and R.~E.~Carlson,
\emph{Monotone piecewise cubic interpolation},
SIAM J.\ Numer.\ Anal.\ \textbf{17}, 238--246 (1980).

\bibitem{Durkan2019}
C.~Durkan, A.~Bekasov, I.~Murray, and G.~Papamakarios,
\emph{Neural spline flows},
Advances in Neural Information Processing Systems (NeurIPS) \textbf{32} (2019).

\bibitem{Papamakarios2021}
G.~Papamakarios, E.~Nalisnick, D.~J.~Rezende, S.~Mohamed, and
B.~Lakshminarayanan,
\emph{Normalizing flows for probabilistic modeling and inference},
J.\ Mach.\ Learn.\ Res.\ \textbf{22}, 1--64 (2021).

\end{thebibliography}

\end{document}
